% =============================================================
%  ch05a_regelungstechnik.tex  --  Regelungstechnik für Roboter:
%  Feedback, Vorsteuerung, Kaskade, Arbeitsraum, Feedback-Entkopplung,
%  Beobachter und Störgrößenkompensation
% =============================================================

\chapter{Regelungstechnik der Robotergrundlagen}
\label{ch:regelung}

Die Dynamik aus Kapitel~\ref{ch:kinematik} sagt, \emph{welches} Gelenkmoment
$\bm{\tau}$ eine gewünschte Bewegung erzeugt. Die Regelung beantwortet die
schwierigere Frage: Wie hält der Roboter die Sollbahn \textbf{trotz}
Modellfehlern, Reibung, Nutzlastwechseln und Stößen? Dieses Kapitel ordnet die
gängigen Regelungsarchitekturen ein -- mit dem Fokus, \textbf{wann} man welche
einsetzen \emph{muss} und worin der Unterschied zu den Alternativen liegt.

Grundlage ist überall das Bewegungsgesetz aus dem Kinematik-Kapitel:
\begin{equation}\label{eq:dyn}
\bm{M}(\bm{q})\ddot{\bm{q}} + \bm{C}(\bm{q},\dot{\bm{q}})\dot{\bm{q}}
+ \bm{g}(\bm{q}) = \bm{\tau} + \bm{\tau}_{\text{stör}}
\end{equation}
mit Massenmatrix $\bm{M}$, Coriolis/Zentrifugal-Term $\bm{C}$, Gravitation
$\bm{g}$ und einer Störung $\bm{\tau}_{\text{stör}}$ (Reibung, Kontaktkräfte,
unmodellierte Last). Der Bahnfehler ist $\bm{e} = \bm{q}_{\text{soll}} - \bm{q}$.

\section{Feedback: die Grundlage}
Feedback (Rückführung) misst den Fehler $\bm{e}$ und stellt ihm ein
korrigierendes Moment entgegen. Der Klassiker im Gelenkraum ist ein
\textbf{PD-Regler mit Gravitationskompensation}:
\begin{equation}
\bm{\tau} = \bm{K}_p\,\bm{e} + \bm{K}_d\,\dot{\bm{e}} + \bm{g}(\bm{q})
\end{equation}
\begin{infobox}[title=Wann einsetzen]
Immer -- Feedback ist das Fundament jeder Roboterregelung. Es ist die einzige
Komponente, die \emph{unbekannte} Störungen und Modellfehler nachträglich
ausbügelt. Ein reiner PD/PID im Gelenkraum genügt bei langsamen Bewegungen,
hoher Getriebeuntersetzung (entkoppelt die Gelenke quasi voneinander) und
moderaten Genauigkeitsanforderungen -- etwa beim Anfahren von Wegpunkten.
\end{infobox}
\textbf{Unterschied zu den anderen:} Feedback reagiert \emph{erst nach} dem
Auftreten des Fehlers -- es ist reaktiv und damit zwangsläufig
„hinterher``. Hohe $\bm{K}_p$ verkürzen die Reaktion, fachen aber Rauschen und
Schwingungen an. Alles Folgende (Vorsteuerung, Entkopplung, Beobachter) dient
dazu, dem Feedback weniger Arbeit zu hinterlassen.

\section{Vorsteuerung (Feedforward)}
Die Vorsteuerung rechnet aus der \emph{bekannten} Sollbahn das nötige Moment
\textbf{vorab} aus und schaltet es additiv zum Feedback. Mit dem Modell
\eqref{eq:dyn} ergibt sich die ideale \textbf{dynamische Vorsteuerung}:
\begin{equation}\label{eq:ff}
\bm{\tau} = \underbrace{\bm{M}(\bm{q})\ddot{\bm{q}}_{\text{soll}}
+ \bm{C}(\bm{q},\dot{\bm{q}})\dot{\bm{q}}_{\text{soll}} + \bm{g}(\bm{q})}
_{\text{Vorsteuerung (Modell)}}
+ \underbrace{\bm{K}_p\,\bm{e} + \bm{K}_d\,\dot{\bm{e}}}_{\text{Feedback}}
\end{equation}
Das ist eine \textbf{Zwei-Freiheitsgrade-Struktur}: Die Vorsteuerung bestimmt
das Führungsverhalten (Bahnfolge), das Feedback nur noch das Störverhalten
(Restfehler). Beide lassen sich getrennt auslegen.
\begin{infobox}[title=Wann einsetzen]
Sobald die Bahn \emph{im Voraus bekannt} ist und schnell/präzise gefolgt werden
soll -- Schweißnähte, Pick-and-Place mit kurzen Taktzeiten, der Schlagarm des
Tischtennisroboters. Die Vorsteuerung eliminiert den \emph{systematischen}
Schleppfehler, ohne die Feedback-Verstärkung hochziehen zu müssen.
\end{infobox}
\textbf{Unterschied zum Feedback:} Vorsteuerung ist \emph{vorausschauend} und
kennt keinen Messfehler -- sie kann daher nicht instabil werden, aber auch keine
unvorhergesehene Störung ausgleichen. Sie ist nur so gut wie das Modell. Darum
nie allein, sondern stets \emph{kombiniert} mit Feedback.

\section{Kaskadenregelung}
Statt einer einzigen Schleife schachtelt die Kaskade mehrere ineinander -- jede
mit eigener, messbarer Zwischengröße. Beim Servoantrieb sind das drei Ebenen:
\begin{center}
\begin{tikzpicture}[node distance=7mm]
\node[node] (pos) {Positions-\\regler};
\node[node,right=of pos] (vel) {Geschwindigkeits-\\regler};
\node[node,right=of vel] (cur) {Strom-/Momenten-\\regler};
\node[node,right=of cur,fill=accent!18] (mot) {Motor +\\Gelenk};
\draw[flow] (pos)--(vel) node[midway,above,font=\scriptsize]{$\dot q_{\text{soll}}$};
\draw[flow] (vel)--(cur) node[midway,above,font=\scriptsize]{$i_{\text{soll}}$};
\draw[flow] (cur)--(mot) node[midway,above,font=\scriptsize]{$u$};
\draw[flow] (mot.south) -- ++(0,-6mm) -| (vel.south)
  node[near start,below,font=\scriptsize]{$\dot q$ (Encoder)};
\draw[flow] (mot.south) -- ++(0,-11mm) -| (pos.south)
  node[near start,below,font=\scriptsize]{$q$ (Encoder)};
\end{tikzpicture}
\end{center}
Die \textbf{innere} Schleife (Strom) ist am schnellsten getaktet
($\sim$\,kHz--10\,kHz), die \textbf{äußere} (Position) am langsamsten. Faustregel:
jede innere Schleife mindestens 3--5$\times$ schneller als die nächstäußere.
\begin{infobox}[title=Wann einsetzen]
Praktisch in jedem industriellen Antrieb -- Servoverstärker und FOC-Treiber sind
intern bereits kaskadiert. Sinnvoll, sobald eine schnell messbare Zwischengröße
(Strom, Drehzahl) existiert und die Stellgröße begrenzt werden muss.
\end{infobox}
\textbf{Unterschied zu den anderen:} Die Kaskade ist eine \emph{strukturelle}
Anordnung mehrerer Feedback-Schleifen, kein eigenständiges Regelgesetz. Ihr
Vorteil: Störungen, die in einer inneren Schleife angreifen (z.\,B.
Versorgungsspannungs\-schwankungen), werden dort sofort abgefangen, bevor sie die
äußere erreichen; außerdem lässt sich jede Ebene einzeln in Betrieb nehmen und
sättigen (Anti-Windup, Strombegrenzung). Quer\-kopplungen zwischen den Gelenken
behandelt sie aber nicht -- dafür braucht es Entkopplung (s.\,u.).

\section{Feedback-Entkopplung (Computed Torque)}
Bei schnellen Bewegungen sind die Gelenke über $\bm{M}$ und $\bm{C}$ stark
\emph{gekoppelt}: ein bewegtes Gelenk übt Reaktionsmomente auf die anderen aus.
Die Feedback-Entkopplung (auch \textbf{inverse Dynamik}- oder
\textbf{Computed-Torque}-Regelung) hebt das per exakter Modellrechnung auf. Man
setzt eine Hilfsgröße $\bm{a}$ an und invertiert die Dynamik:
\begin{equation}\label{eq:ctc}
\bm{\tau} = \bm{M}(\bm{q})\,\bm{a}
+ \bm{C}(\bm{q},\dot{\bm{q}})\dot{\bm{q}} + \bm{g}(\bm{q}),
\qquad
\bm{a} = \ddot{\bm{q}}_{\text{soll}} + \bm{K}_d\,\dot{\bm{e}} + \bm{K}_p\,\bm{e}
\end{equation}
Eingesetzt in \eqref{eq:dyn} (ohne Störung) bleibt das exakt \textbf{lineare,
entkoppelte} Fehlersystem
\begin{equation}
\ddot{\bm{e}} + \bm{K}_d\,\dot{\bm{e}} + \bm{K}_p\,\bm{e} = \bm{0},
\end{equation}
also für \emph{jedes} Gelenk ein unabhängiges, frei platzierbares
PD-Verhalten -- unabhängig von Stellung und Geschwindigkeit.
\begin{infobox}[title=Wann einsetzen]
Schnelle, dynamische Bewegungen mit leichter Struktur und geringer
Getriebeunter\-setzung, bei denen die Kopplung spürbar wird -- Leichtbauarme,
Direktantriebe, der hochbeschleunigte Schlagarm. Voraussetzung: ein hinreichend
genaues Dynamikmodell und Momenten-/Strom\-regelbare Antriebe.
\end{infobox}
\textbf{Unterschied zur Vorsteuerung:} Beide nutzen das Modell. Die Vorsteuerung
\eqref{eq:ff} wertet es entlang der \emph{Soll}-Größen aus (Steuerung, offene
Wirkung); die Feedback-Entkopplung \eqref{eq:ctc} wertet $\bm{M},\bm{C}$ entlang
der \emph{Ist}-Größen aus und linearisiert die Strecke \emph{innerhalb} der
geschlossenen Schleife. Dadurch ist das resultierende Fehlerverhalten
arbeitspunktunabhängig konstant -- der Preis ist die Modellauswertung in
Echtzeit und die Empfindlichkeit gegen Modellfehler.

\section{Regelung im Arbeitsraum (Operational Space Control)}
Bisher wurde im Gelenkraum geregelt; die Aufgabe ist aber meist im
\textbf{kartesischen Raum} definiert (Pose des Endeffektors $\bm{x}$). Statt erst
per inverser Kinematik nach $\bm{q}_{\text{soll}}$ umzurechnen, formuliert man
Regler und Dynamik direkt in $\bm{x}$ -- über die Jacobi-Matrix $\bm{J}$ mit
$\dot{\bm{x}} = \bm{J}\dot{\bm{q}}$. Das entkoppelte Gesetz im Arbeitsraum lautet
\begin{equation}
\bm{\tau} = \bm{J}^\top\!\Big[\bm{\Lambda}(\bm{q})\,\bm{a}_x
+ \bm{\mu}(\bm{q},\dot{\bm{q}}) + \bm{p}(\bm{q})\Big],
\qquad
\bm{a}_x = \ddot{\bm{x}}_{\text{soll}}
+ \bm{K}_d\,\dot{\bm{e}}_x + \bm{K}_p\,\bm{e}_x
\end{equation}
mit der Arbeitsraum-Massenmatrix $\bm{\Lambda} = (\bm{J}\bm{M}^{-1}\bm{J}^\top)^{-1}$.
Der Fehler $\bm{e}_x$ wird direkt im kartesischen Raum gebildet (Position als
Differenz, Orientierung als Vektorteil des Fehlerquaternions, vgl.
Kapitel~\ref{ch:kinematik}).
\begin{infobox}[title=Wann einsetzen]
Wenn die \emph{Aufgabe} im kartesischen Raum lebt: Kraft-/Impedanz\-regelung beim
Kontakt (Polieren, Stecken, Montage), Bahnverfolgung am Werkstück,
nullraumoptimierte redundante Arme. Auch dann, wenn die inverse Kinematik nahe
Singularitäten unzuverlässig wird.
\end{infobox}
\textbf{Unterschied zur Gelenkraumregelung:} Identische Mathematik, anderer Raum.
Im Arbeitsraum legst du Steifigkeit und Dämpfung \emph{richtungsabhängig} fest
(z.\,B.\ nachgiebig in Stoßrichtung, steif quer dazu) -- ideal für Kontakt\-aufgaben.
Nachteil: $\bm{J}$ wird nahe Singularitäten schlecht konditioniert, und der
Nullraum redundanter Roboter muss separat behandelt werden.

\section{Zustandsbeobachter}
Computed-Torque und Vorsteuerung brauchen $\dot{\bm{q}}$ (oft auch
Beschleunigungen oder die Störung) -- gemessen wird aber meist nur die Position
$\bm{q}$. Numerisches Differenzieren verstärkt Rauschen. Der
\textbf{Zustandsbeobachter} schätzt die nicht gemessenen Zustände aus Modell und
Eingang, korrigiert über den Messfehler:
\begin{equation}
\dot{\hat{\bm{x}}} = \bm{f}(\hat{\bm{x}},\bm{u})
+ \bm{L}\,(\bm{y} - \hat{\bm{y}})
\end{equation}
Die Beobachterverstärkung $\bm{L}$ legt fest, wie schnell die Schätzung dem
wahren Zustand folgt (Luenberger für lineare, EKF für nichtlineare Systeme --
siehe Kapitel~\ref{ch:kalman}).
\begin{infobox}[title=Wann einsetzen]
Wenn Zustände für das Regelgesetz gebraucht werden, die kein Sensor direkt
liefert: Gelenkgeschwindigkeit aus reinem Positionsencoder, Verspannung in einem
elastischen Gelenk, oder -- besonders wichtig -- die \emph{Störung} selbst (s.\,u.).
\end{infobox}
\textbf{Unterschied zum Kalman-Filter:} Der Beobachter ist das deterministische
Konzept; der Kalman-Filter ist der Beobachter, dessen $\bm{L}$ statistisch
optimal aus den Rausch-Kovarianzen bestimmt wird. Funktional dasselbe Bauteil im
Regelkreis.

\section{Störgrößenkompensation}
Reibung, unbekannte Nutzlast und Kontaktkräfte stecken in
$\bm{\tau}_{\text{stör}}$ aus \eqref{eq:dyn}. Drei Strategien -- gestaffelt nach
dem, was über die Störung bekannt ist:

\subsection{Statische vs.\ dynamische Störgrößenausschaltung}
Lässt sich die Störung \emph{messen} oder aus den Eingängen \emph{berechnen}
(z.\,B.\ ein bekanntes Lastmoment, das von der nachgeführten Last abhängt),
schaltet man ihren Effekt vorausberechnet auf -- analog zur Vorsteuerung, nur für
die Störung statt für die Führung. Die \textbf{dynamische} Variante berücksichtigt
dabei das \emph{zeitliche} Übertragungsverhalten zwischen Störangriff und Ausgang
(nicht nur den stationären Wert), sodass die Kompensation auch während des
Einschwingens passt.
\begin{infobox}[title=Wann einsetzen]
Wenn die Störung \emph{messbar} ist (Kraftsensor, bekannte, mitgeführte Last) und
schnell genug wirkt, dass reines Feedback zu spät käme. Die statische Form genügt
bei langsamen Störungen; die dynamische, wenn Störung und Strecke vergleichbar
schnell sind.
\end{infobox}

\subsection{Störentkopplung (Disturbance Decoupling)}
Hier ist das Ziel strenger: die Reglerstruktur so wählen, dass die Störung den
\emph{geregelten Ausgang} \textbf{prinzipiell nicht mehr beeinflusst} -- die
Übertragungsfunktion von $\bm{\tau}_{\text{stör}}$ auf $\bm{y}$ wird zu null. Das
gelingt nur, wenn der Störangriff aus dem messbaren/stellbaren Raum heraus
kompensierbar ist (die geometrische Entkopplungsbedingung). Die Störung darf
intern noch wirken, erscheint aber nicht am Ausgang.
\begin{infobox}[title=Wann einsetzen]
Wenn eine bestimmte Störrichtung den Ausgang \emph{strukturell} nicht erreichen
darf und das System die Entkopplungsbedingung erfüllt. Schärfer als die bloße
Ausschaltung, dafür an strukturelle Voraussetzungen gebunden.
\end{infobox}
\textbf{Unterschied:} Ausschaltung \emph{kompensiert} eine konkrete Störung
betragsmäßig; Entkopplung \emph{eliminiert} ihren Einfluss auf den Ausgang per
Struktur -- unabhängig vom Störverlauf.

\subsection{Dynamische Vorsteuerung -- mit oder ohne Zustandsregelung}
Ist die Störung \emph{nicht} messbar, schätzt sie ein Beobachter (oft als
\textbf{Disturbance Observer}, DOB, der $\hat{\bm{\tau}}_{\text{stör}}$ aus dem
Vergleich von Modell und Istverhalten rekonstruiert) und speist die Schätzung
gegenphasig auf:
\begin{equation}
\bm{\tau} = \bm{\tau}_{\text{Regler}} - \hat{\bm{\tau}}_{\text{stör}}.
\end{equation}
Das ist eine \emph{dynamische Vorsteuerung der Störung} aus dem Beobachter.

Zwei Ausbaustufen:
\begin{itemize}
\item \textbf{Ohne Zustandsregelung:} reine Aufschaltung der geschätzten Störung
  auf einen bestehenden (z.\,B.\ PD-) Regler. Einfach nachzurüsten, sehr wirksam
  gegen Reibung und langsame Lastwechsel.
\item \textbf{Mit Zustandsregelung:} die Störung wird als \emph{erweiterter
  Zustand} ins Modell aufgenommen (Störmodell), der Beobachter schätzt sie
  mit, und ein Zustandsregler $\bm{u} = -\bm{K}\hat{\bm{x}}$ verarbeitet sie
  konsistent mit den übrigen Zuständen. Das erlaubt frei platzierbare Dynamik des
  Gesamtsystems inklusive Störunterdrückung.
\end{itemize}
\begin{warnbox}[title=Voraussetzung und Grenze]
Jede modellbasierte Störschätzung lebt vom Modell: Ist $\bm{M},\bm{C},\bm{g}$
falsch, interpretiert der DOB den Modellfehler als Störung und „kompensiert`` ihn
ebenfalls -- gewollt bei Robustheit, ungewollt, wenn die Bandbreite zu hoch
gewählt ist und Messrauschen mit aufgeschaltet wird.
\end{warnbox}
\textbf{Unterschied zur messenden Ausschaltung:} Statt die Störung zu
\emph{messen}, wird sie \emph{geschätzt} -- anwendbar auch dort, wo kein Sensor
sie erfasst (Reibung!). Die Zustandsregelungs-Variante unterscheidet sich von der
reinen Aufschaltung dadurch, dass Stör- und Regeldynamik \emph{gemeinsam}
ausgelegt werden, statt eine fertige Schleife nur zu ergänzen.

\section{Überblick: Wann was?}
\begin{longtable}{@{}p{3.4cm}p{5.3cm}p{5.6cm}@{}}
\toprule
\textbf{Verfahren} & \textbf{Einsatz / Auslöser} & \textbf{Abgrenzung} \\
\midrule
Feedback (PD/PID) & Immer; langsame Bewegung, hohe Untersetzung &
  Reaktiv, gleicht Unbekanntes aus, aber „hinterher`` \\
Vorsteuerung & Bahn vorab bekannt, schnelle Folge nötig &
  Vorausschauend, modellbasiert, ohne Feedback nicht störfest \\
Kaskade & Schnelle Zwischengröße messbar (Strom, Drehzahl) &
  Struktur aus Schleifen, fängt innere Störungen früh ab \\
Feedback-Entkopplung & Schnelle Bewegung, spürbare Gelenk\-kopplung &
  Linearisiert Strecke \emph{im} Kreis, arbeitspunktunabhängig \\
Arbeitsraum-Regelung & Aufgabe kartesisch, Kontakt/Kraft &
  Wie Entkopplung, aber in $\bm{x}$; richtungsabhängige Steifigkeit \\
Zustandsbeobachter & Zustand/Störung nicht direkt messbar &
  Liefert $\dot{\bm{q}}$, Störung etc.; KF = optimaler Beobachter \\
Störausschaltung & Störung messbar &
  Kompensiert Betrag (statisch/dynamisch) \\
Störentkopplung & Störrichtung darf Ausgang nicht erreichen &
  Eliminiert Einfluss strukturell, nicht nur betragsmäßig \\
Dyn.\ Vorst.\ (DOB) & Störung nicht messbar (Reibung, Last) &
  Schätzt statt misst; mit Zustandsregelung gemeinsam ausgelegt \\
\bottomrule
\end{longtable}

\noindent\textbf{Praxis-Fazit:} Die Verfahren sind \emph{Schichten}, keine
Alternativen. Ein realer Roboterregler kombiniert sie: kaskadierte
Antriebsregelung als Basis, Feedback-Entkopplung oder Arbeitsraumregelung für
die Kopplung, Vorsteuerung für die bekannte Bahn, ein Beobachter für fehlende
Zustände und ein Disturbance Observer für Reibung und Last. Jede Schicht nimmt
dem darunterliegenden Feedback Arbeit ab -- und nur das Feedback bleibt am Ende
für das Unvorhergesehene zuständig.
